\documentclass[14pt]{extarticle}
\usepackage{fontspec}
\usepackage{geometry}
\geometry{a4paper, margin=2.0cm}
\usepackage{graphicx}
\usepackage{booktabs}
\usepackage{array}
\usepackage{longtable}
\usepackage{amssymb}
\usepackage{amsmath}
\usepackage{fancyvrb}
\usepackage[table]{xcolor}
\usepackage{titlesec}
\usepackage{fancyhdr}
\usepackage{enumitem}
\usepackage[most]{tcolorbox}
\usepackage{hyperref}

% ---- Cool theme palette (calm teal / slate — not "hot") ----
\definecolor{theme}{HTML}{2b6a77}    % primary teal-slate
\definecolor{themeD}{HTML}{1f4e58}   % darker teal
\definecolor{ink}{HTML}{232d33}      % near-black slate
\definecolor{ok}{HTML}{3f8161}       % muted green (done)
\definecolor{cardbg}{HTML}{f3f6f7}   % very light cool gray
\definecolor{softgray}{HTML}{eef1f2}
\definecolor{rowtint}{HTML}{e8eef0}

\hypersetup{colorlinks=true, linkcolor=theme, urlcolor=theme}

% ---- Fonts ----
\setmainfont{THSarabunNew.ttf}[
  Path=./fonts/, BoldFont=THSarabunNew Bold.ttf,
  ItalicFont=THSarabunNew Italic.ttf, BoldItalicFont=THSarabunNew BoldItalic.ttf]
\setmonofont{cour.ttf}[Path=./fonts/, BoldFont=courbd.ttf, Scale=0.80]
\XeTeXlinebreaklocale "th"
\XeTeXlinebreakskip = 0pt plus 0.1pt

\titleformat{\section}{\Large\bfseries\color{theme}}{\thesection.}{0.5em}{}
\titleformat{\subsection}{\large\bfseries\color{ink}}{\thesubsection}{0.5em}{}
\setlist{itemsep=1pt, topsep=2pt}
\pagestyle{fancy}\fancyhf{}
\fancyhead[L]{\small\color{ink} SPR Fracture Detection — Experimental Checklist}
\fancyhead[R]{\small\color{ink}\thepage}
\renewcommand{\headrulewidth}{0.4pt}
\renewcommand{\headrule}{\hbox to\headwidth{\color{theme!45}\leaders\hrule height \headrulewidth\hfill}}

\newcommand{\todo}{\textcolor{theme}{$\square$}}
\newcommand{\done}{\textcolor{ok}{$\boxtimes$}}
% \code — literal mono (underscores etc. need no escaping)
\newcommand{\code}[1]{{\ttfamily\detokenize{#1}}}

% Verbatim run-command style
\DefineVerbatimEnvironment{run}{Verbatim}{frame=single, framerule=0.4pt,
  rulecolor=\color{theme!40}, fontsize=\small, formatcom=\color{ink}, xleftmargin=6pt}

% Experiment card
\newtcolorbox{card}[1]{colback=cardbg, colframe=theme!55, boxrule=0.8pt, arc=2pt,
  left=8pt, right=8pt, top=5pt, bottom=6pt, title=#1, coltitle=white,
  fonttitle=\bfseries, colbacktitle=theme}

% sub-run table helper column types
\newcolumntype{K}{>{\centering\arraybackslash}p{0.7em}}

\begin{document}

% ===================== TITLE =====================
\begin{center}
{\LARGE\bfseries\color{theme} ระเบียบวิธีการทดลองเพิ่มเติม (Experimental Checklist)}\\[4pt]
{\Large การตรวจจับรอยแตกกระดูกเชิงกรานด้วยลายเซ็นการกระเจิงรังสีเอกซ์ (SPR)}\\[6pt]
{\small แผนการทดลองสำหรับผู้ร่วมวิจัย --- ใช้เฉพาะโมเดลที่มีอยู่ปัจจุบัน โดยปรับพารามิเตอร์การรัน}\\
{\small ระบบจำลอง OpenGATE 10 / Geant4 \quad$\bullet$\quad \today}
\end{center}
\vspace{2pt}{\color{theme!45}\hrule}\vspace{8pt}

% ===================== SECTION 0 =====================
\section{บทนำและวิธีใช้เอกสาร}
เอกสารนี้เป็น \textbf{check list ของการทดลอง} เพื่อทดสอบสมมติฐานหลัก---``รอยแตกกระดูกทิ้งลายเซ็นที่
ตรวจจับได้ในอัตราส่วนการกระเจิงต่อรังสีปฐมภูมิ (SPR)''---ให้\textbf{ครอบคลุมและแม่นยำยิ่งขึ้น}
ภายใต้เงื่อนไข\textbf{ใช้เพียงโมเดลที่มีอยู่} (รอยแตก 4 ขนาด + matched control) แล้ว\textbf{ปรับพารามิเตอร์การรัน}

\noindent\textbf{วิธีใช้:} แต่ละการทดลองใหญ่ (CL-x) มี \emph{สมมติฐาน, ทฤษฎี, ผลที่ทำนาย} แล้วแตกเป็น
\textbf{รายการย่อย (CL-x.n) หนึ่งแถวต่อหนึ่งการรัน} พร้อมพารามิเตอร์ที่ต้องตั้ง --- ผู้ร่วมวิจัยรัน
ตามลำดับความสำคัญ (\S\ref{sec:priority}) ติ๊ก \todo\,$\to$\,\done\ ทีละแถว และบันทึกผลลงตาราง (\S\ref{sec:log})

\noindent\textbf{สัญลักษณ์:} \done~= ทำแล้ว \quad \todo~= รอรัน

\begin{tcolorbox}[colback=softgray, colframe=theme!25, boxrule=0.5pt, arc=2pt]
\textbf{หลักการร่วมของทุกการทดลอง (ห้ามละ):} (1)~\textbf{matched control} (ต่างจาก fracture เฉพาะร่อง)
ตัด confound ของ mesh; (2)~\textbf{multi-seed} (control+fracture ใช้ seed เดียวกันต่อ realization =
common random numbers) วัดนัยสำคัญ; (3)~\textbf{a-priori ROI} ที่ตำแหน่งร่อง กำหนดก่อนดูผล เทียบ background ---
ห้ามเลือก ROI จากจุดสัญญาณแรงสุด (cherry-picking)
\end{tcolorbox}

% ===================== SECTION 1: THEORY =====================
\section{ทฤษฎีพื้นฐาน}

\subsection{อัตราส่วนการกระเจิง SPR และกลไกของรอยแตก}
\[ \mathrm{SPR}(x,y)=\frac{I_{\text{scatter}}(x,y)}{I_{\text{primary}}(x,y)} \]
รอยแตก = ช่องอากาศบางในกระดูก; Beer--Lambert $I=I_0 e^{-\int \mu\,dl}$ โดย
$\mu_{\text{air}}\!\ll\!\mu_{\text{bone}}$ $\Rightarrow$ ตามรังสีที่ผ่านร่อง primary เพิ่ม $\Rightarrow$
$\Delta\mathrm{SPR}=\mathrm{SPR}_{\text{frac}}-\mathrm{SPR}_{\text{ctrl}}<0$ ที่ร่อง (ตรงกับ CL-0)

\subsection{การพึ่งพาพลังงาน}
$\mu_{\text{bone}}(E)=\underbrace{\tau}_{\text{photoelec}}+\underbrace{\sigma_{\text{C}}}_{\text{Compton}}+\dots$,
$\tau\propto Z^{\sim4}/E^{\sim3}$ (ตกเร็วเมื่อ $E$ สูง) $\Rightarrow$ contrast ร่องแรงที่ $E$ ต่ำ,
scatter fraction (SPR ฐาน) เพิ่มที่ $E$ สูง $\Rightarrow$ มีพลังงานเหมาะสมสำหรับ CNR (CL-3)

\subsection{การกระเจิงกับเรขาคณิต}
scatter $\propto$ ปริมาตรฉาย $\Rightarrow$ \textbf{collimation} ลด scatter บดบัง (CL-4);
\textbf{air-gap/ODD} ให้ scatter ที่ฟุ้งพลาดฉากรับ (CL-5);
\textbf{มุมลำแสง} --- สัญญาณ $\propto$ เส้นทางอากาศที่ฉายผ่านร่อง สูงสุดเมื่อลำแสงเลียบระนาบร่อง (CL-6)

\subsection{กรอบสถิติ}
\[ t=\frac{\overline{\Delta\mathrm{SPR}}}{s/\sqrt{N}},\quad \mathrm{df}=N-1;\qquad
   \text{s.e.}\propto\tfrac{1}{\sqrt{N_{\text{seeds}}}},\quad
   \text{MC noise}\propto\tfrac{1}{\sqrt{N_{\text{photons}}}} \]
เพิ่ม seed/photons $\Rightarrow$ $|t|$ โตถ้าสัญญาณจริง (CL-1); control-vs-control/permutation คุม false-positive (CL-2)

% ===================== SECTION 2: MASTER TABLE =====================
\section{ตารางสรุป (Master Checklist)}
\renewcommand{\arraystretch}{1.2}
\small
\begin{longtable}{c l l c c}
\rowcolor{theme}\color{white}\textbf{ID} & \color{white}\textbf{การทดลอง} &
  \color{white}\textbf{ทดสอบอะไร} & \color{white}\textbf{\#runs} & \color{white}\textbf{สถานะ} \\
\midrule
\endhead
CL-0 & Baseline detectability sweep & สัญญาณจริง \& โตตามขนาด & 4 & \done \\
\rowcolor{rowtint}
CL-1 & Statistical convergence & ยืนยัน 0.306 \& คลี่คลาย plateau & 4 & \todo \\
CL-2 & Null / false-positive & อัตราผลบวกลวงต่ำ (validity) & 3 & \todo \\
\rowcolor{rowtint}
CL-3 & Energy dependence & กลไกฟิสิกส์ \& พลังงานเหมาะสม & 5 & \todo \\
CL-4 & Collimation field-size & collimation ช่วยจริง \& ค่าเหมาะสม & 6 & \todo \\
\rowcolor{rowtint}
CL-5 & Air-gap / ODD & scatter rejection เชิงเรขาคณิต & 5 & \todo \\
CL-6 & Beam-angle / orientation & \textbf{ทดสอบเชิงสาเหตุ}ทิศทาง & 7 & \todo \\
\rowcolor{rowtint}
CL-7 & ROI \& spatial profile & matched filter \& อธิบาย plateau & 3 & \todo \\
CL-8 & Detector pixel-size & partial-volume \& ความละเอียด & 4 & \todo \\
\bottomrule
\end{longtable}
\normalsize

% ===================== SECTION 3: CARDS =====================
\section{รายละเอียดการทดลอง (พร้อมรายการย่อย)}

% ---------- CL-0 ----------
\begin{card}{CL-0 \quad Baseline matched-control sweep \hfill \done\ ทำแล้ว}
\textbf{สมมติฐาน:} สัญญาณ $\Delta\mathrm{SPR}$ ตรวจพบได้อย่างมีนัยสำคัญและโตตามขนาดร่อง \quad
\textbf{Baseline:} \code{energy_keV=80, photons=2000000, pix=256, field_mm=40, n_seeds=8, sod=800, odd=400}
\end{card}
\vspace{-1pt}\renewcommand{\arraystretch}{1.15}\small
\begin{center}
\begin{tabular}{cc c c c c}
\rowcolor{softgray}\ & ขนาดร่อง & $\Delta$SPR & $|t|$ & $p$ & bg \\
\done & 0.306 & $-3.3\times10^{-5}$ & 3.12 & 0.017 & null \\
\done & 0.43  & $-3.4\times10^{-5}$ & 3.66 & 0.008 & null \\
\done & 0.54  & $-4.9\times10^{-5}$ & 4.84 & 0.002 & null \\
\done & 1.0   & $-4.2\times10^{-5}$ & 4.32 & 0.004 & null \\
\end{tabular}
\end{center}
\normalsize
\textbf{ค้นพบ:} ตรวจพบครบ ($p<0.02$), เครื่องหมายลบตามฟิสิกส์, background null, อิ่มตัวที่ $\ge0.54$mm,
ตัวชี้วัดพื้นผิว (variance/entropy) บอด (สัญญาณเป็น mean-shift) $\to$ นำสู่ CL-1 \& CL-7

\vspace{6pt}
% ---------- CL-1 ----------
\begin{card}{CL-1 \quad Statistical convergence \& power \hfill \todo}
\textbf{สมมติฐาน:} เพิ่มสถิติ $\Rightarrow$ $|t|$ โตเสถียร ยืนยัน 0.306mm และชี้ขาด plateau (0.54 vs 1.0)\\
\textbf{ทฤษฎี:} $\text{s.e.}\propto1/\sqrt{N_{\text{seeds}}}$, noise $\propto1/\sqrt{N_{\text{photons}}}$;
$t\propto$ สัญญาณ$\times\sqrt{N}$\\
\textbf{ทำนาย:} 0.306mm $|t|$ เกินเกณฑ์มั่นคง; plateau ชี้ขาด (ถ้า 1.0 ยังไม่เกิน 0.54 = อิ่มตัวจริง $\to$ CL-7)\\
\textbf{ตรึงทุกแถว:} \code{field_mm=40, sod=800, odd=400}
\smallskip\\ \textbf{รายการย่อย:}
\end{card}
\vspace{-1pt}\small
\begin{tabular}{K l l l}
\todo & \textbf{1.1} & โมเดล 0.306 & \code{n_seeds=30, photons=2000000} \\
\todo & \textbf{1.2} & โมเดล 0.306 & \code{n_seeds=30, photons=8000000} \\
\todo & \textbf{1.3} & โมเดล 0.54  & \code{n_seeds=30, photons=8000000} \\
\todo & \textbf{1.4} & โมเดล 1.0   & \code{n_seeds=30, photons=8000000} \\
\end{tabular}
\normalsize

\vspace{6pt}
% ---------- CL-2 ----------
\begin{card}{CL-2 \quad Null / false-positive calibration \hfill \todo\ (ทำก่อน)}
\textbf{สมมติฐาน:} ระเบียบวิธีไม่สร้างสัญญาณลวงเมื่อไม่มีรอยแตก\\
\textbf{ทฤษฎี:} ต้องคุม false-positive rate ก่อนเชื่อผลบวก --- รันแบบ ``ไม่มีสัญญาณ'' ต้องได้ null\\
\textbf{ทำนาย:} 2.1/2.3 ให้ $t\approx0,\ p>0.05$; 2.2 ตำแหน่งร่องเป็น outlier ชัด (empirical $p$)\\
\textbf{รายการย่อย:} (2.1--2.2 = วิเคราะห์ข้อมูล CL-0 เดิม ไม่ต้องรันใหม่)
\end{card}
\vspace{-1pt}\small
\begin{tabular}{K l p{10.2cm}}
\todo & \textbf{2.1} & \textbf{permutation-ROI} (วิเคราะห์): วาง ROI 16$\times$16 สุ่ม $\ge200$ ตำแหน่งในลำแสง (เลี่ยงร่อง) บนข้อมูล CL-0 $\to$ null distribution ของ $t$ \\
\todo & \textbf{2.2} & \textbf{split-seed} (วิเคราะห์): แยก 8 seeds ของ control เป็น 2 กลุ่ม หาผลต่างระหว่างกลุ่ม $\to$ ต้อง null \\
\todo & \textbf{2.3} & \textbf{control-vs-control} (รันใหม่, เลือกได้): รัน control model เป็นทั้งสองฝั่งด้วย \code{--random_seed} คนละชุด (offset ผ่าน CLI) $\to$ $\Delta$SPR ควรเป็น noise \\
\end{tabular}
\normalsize

\vspace{6pt}
% ---------- CL-3 ----------
\begin{card}{CL-3 \quad Energy dependence (spectral) \hfill \todo\ (กลไกฟิสิกส์)}
\textbf{สมมติฐาน:} ถ้าสัญญาณจากการลดทอนของกระดูกจริง $|\Delta\mathrm{SPR}|$/CNR ต้องแปรตามพลังงาน\\
\textbf{ทฤษฎี:} $\tau\propto1/E^{\sim3}$ (contrast แรงที่ $E$ ต่ำ), scatter เพิ่มที่ $E$ สูง $\to$ มีจุดเหมาะสม\\
\textbf{ทำนาย:} $|\Delta\mathrm{SPR}|$ ลดเมื่อ $E$ สูง, CNR อาจพีคกลาง --- ตรงฟิสิกส์ = สนับสนุนกลไกแข็งแรง\\
\textbf{ตรึงทุกแถว:} \code{model=0.54, n_seeds=10, field_mm=40, pix=256, sod=800, odd=400}
\smallskip\\ \textbf{รายการย่อย (กวาด \code{energy_keV}):}
\end{card}
\vspace{-1pt}\small
\begin{tabular}{K l l @{\hspace{2em}} K l l}
\todo & \textbf{3.1} & \code{energy_keV=40}  &
\todo & \textbf{3.4} & \code{energy_keV=100} \\
\todo & \textbf{3.2} & \code{energy_keV=60}  &
\todo & \textbf{3.5} & \code{energy_keV=120} \\
\todo & \textbf{3.3} & \code{energy_keV=80}  & & & \\
\end{tabular}
\normalsize

\vspace{6pt}
% ---------- CL-4 ----------
\begin{card}{CL-4 \quad Collimation field-size dependence \hfill \todo}
\textbf{สมมติฐาน:} collimation เพิ่ม detectability จริง โดยมีขนาดสนามเหมาะสม\\
\textbf{ทฤษฎี:} scatter $\propto$ ปริมาตรฉาย --- สนามกว้าง scatter บดบัง contrast; สนามแคบสะอาดแต่สถิติน้อย\\
\textbf{ทำนาย:} $t$/CNR พีคที่สนามแคบ--กลาง แล้วตกเมื่อกว้าง (scatter wash)\\
\textbf{ตรึงทุกแถว:} \code{model=0.54, energy_keV=80, n_seeds=10, pix=256, sod=800, odd=400}
\smallskip\\ \textbf{รายการย่อย (กวาด \code{field_mm}):}
\end{card}
\vspace{-1pt}\small
\begin{tabular}{K l l @{\hspace{2em}} K l l}
\todo & \textbf{4.1} & \code{field_mm=15} &
\todo & \textbf{4.4} & \code{field_mm=60} \\
\todo & \textbf{4.2} & \code{field_mm=25} &
\todo & \textbf{4.5} & \code{field_mm=100} \\
\todo & \textbf{4.3} & \code{field_mm=40} &
\todo & \textbf{4.6} & \code{field_mm=0} (เต็มฟิล์ม) \\
\end{tabular}
\normalsize

\vspace{6pt}
% ---------- CL-5 ----------
\begin{card}{CL-5 \quad Air-gap / object--detector distance (ODD) \hfill \todo}
\textbf{สมมติฐาน:} เพิ่ม ODD ลด scatter ที่ถึงฉากรับ (air-gap) + เพิ่มกำลังขยายร่อง\\
\textbf{ทฤษฎี:} scatter ฟุ้งพลาดฉากรับมากกว่า primary เมื่อ ODD มาก; $m=\mathrm{SID}/\mathrm{SOD}$ เพิ่ม\\
\textbf{ทำนาย:} SPR ฐานลดตาม ODD (ยืนยัน air-gap); CNR อาจดีขึ้นแล้วตก (fluence loss)\\
\textbf{ตรึงทุกแถว:} \code{model=0.54, energy_keV=80, n_seeds=10, field_mm=40, pix=256, sod=800}
\smallskip\\ \textbf{รายการย่อย (กวาด \code{odd}):}
\end{card}
\vspace{-1pt}\small
\begin{tabular}{K l l @{\hspace{2em}} K l l}
\todo & \textbf{5.1} & \code{odd=100} &
\todo & \textbf{5.4} & \code{odd=700} \\
\todo & \textbf{5.2} & \code{odd=250} &
\todo & \textbf{5.5} & \code{odd=1200} \\
\todo & \textbf{5.3} & \code{odd=400} & & & \\
\end{tabular}
\normalsize

\vspace{6pt}
% ---------- CL-6 ----------
\begin{card}{CL-6 \quad Beam-angle / fracture-orientation \hfill \todo\ (เชิงสาเหตุแข็งแรงสุด)}
\textbf{สมมติฐาน:} สัญญาณขึ้นกับมุมระหว่างลำแสงกับระนาบร่อง --- สูงสุดเมื่อเลียบระนาบ (tangential)\\
\textbf{ทฤษฎี:} $\Delta(\text{primary})\propto$ เส้นทางอากาศที่ฉายผ่านร่อง = ฟังก์ชันของมุม; ตรงหลักคลินิก\\
\textbf{ทำนาย:} $|\Delta\mathrm{SPR}|$ แปรตามมุมเป็นระบบ (พีคที่มุม tangential) = ลายนิ้วมือเชิงสาเหตุ\\
\textbf{ตรึงทุกแถว:} \code{model=0.54, energy_keV=80, n_seeds=10, field_mm=40, sod=800, odd=400};
ตั้ง \code{src_x=800*sin(t)}, \code{src_z=-800*cos(t)}; ใช้ \textbf{Preview} จัดร่องกลางลำแสง (\code{obj_dx/obj_dy}) ทุกมุม
\smallskip\\ \textbf{รายการย่อย (กวาดมุม $\theta$):}
\end{card}
\vspace{-1pt}\small
\begin{tabular}{K l l l}
\todo & \textbf{6.1} & $\theta=0^\circ$   & \code{src_x=0, src_z=-800} \\
\todo & \textbf{6.2} & $\theta=+10^\circ$ & \code{src_x=138.9, src_z=-787.8} \\
\todo & \textbf{6.3} & $\theta=+20^\circ$ & \code{src_x=273.6, src_z=-751.8} \\
\todo & \textbf{6.4} & $\theta=+30^\circ$ & \code{src_x=400, src_z=-692.8} \\
\todo & \textbf{6.5} & $\theta=-10^\circ$ & \code{src_x=-138.9, src_z=-787.8} \\
\todo & \textbf{6.6} & $\theta=-20^\circ$ & \code{src_x=-273.6, src_z=-751.8} \\
\todo & \textbf{6.7} & $\theta=-30^\circ$ & \code{src_x=-400, src_z=-692.8} \\
\end{tabular}
\normalsize

\vspace{6pt}
% ---------- CL-7 ----------
\begin{card}{CL-7 \quad ROI scaling \& spatial profile \hfill \todo\ (อธิบาย plateau)}
\textbf{สมมติฐาน:} plateau เกิดจาก ROI คงที่เจือจางสัญญาณของร่องกว้าง; ขนาดเชิงพื้นที่ยังโตตามจริง\\
\textbf{ทฤษฎี:} matched filter (SNR สูงสุดเมื่อ ROI $\approx$ ขอบเขตสัญญาณ); FWHM ของแอ่งแปรตามความกว้างร่องที่ฉาย\\
\textbf{ทำนาย:} ความลึก/สัญญาณรวมโตต่อเนื่องแม้ค่าเฉลี่ยใน ROI คงที่ = plateau เป็นผลของ ROI\\
\textbf{รายการย่อย:} (ทั้งหมดวิเคราะห์บนข้อมูล CL-0/CL-1 ไม่ต้องรันใหม่)
\end{card}
\vspace{-1pt}\small
\begin{tabular}{K l p{10.2cm}}
\todo & \textbf{7.1} & กวาดขนาด ROI $\in\{8,12,16,24,32\}$px บนทั้ง 4 ขนาด $\to$ $t$ vs ROI \\
\todo & \textbf{7.2} & ROI ปรับตามความกว้างร่องที่ฉายแต่ละขนาด (matched filter) \\
\todo & \textbf{7.3} & line profile ตัดขวางร่อง: วัดความลึกแอ่ง + FWHM เทียบ vs ขนาดร่อง 4 ขนาด \\
\end{tabular}
\normalsize

\vspace{6pt}
% ---------- CL-8 ----------
\begin{card}{CL-8 \quad Detector pixel-size / resolution \hfill \todo}
\textbf{สมมติฐาน:} มีขนาดพิกเซลเหมาะสม --- หยาบเจือจาง (partial volume), ละเอียด noise สูง\\
\textbf{ทฤษฎี:} พิกเซลหยาบเฉลี่ยสัญญาณกับกระดูกรอบ ๆ แต่โฟตอน/พิกเซลมาก; ละเอียดแยกร่องแต่ noisy\\
\textbf{ทำนาย:} CNR พีคที่พิกเซล$\approx$ความกว้างร่องที่ฉาย แล้วตกทั้งสองฝั่ง\\
\textbf{ตรึงทุกแถว:} \code{model=0.54, energy_keV=80, n_seeds=10, field_mm=40, sod=800, odd=400} (\code{film_xy=400})
\smallskip\\ \textbf{รายการย่อย (กวาด \code{pix}):}
\end{card}
\vspace{-1pt}\small
\begin{tabular}{K l l l}
\todo & \textbf{8.1} & \code{pix=128}  & (พิกเซล 3.13 มม.) \\
\todo & \textbf{8.2} & \code{pix=256}  & (พิกเซล 1.56 มม.) \\
\todo & \textbf{8.3} & \code{pix=512}  & (พิกเซล 0.78 มม.) \\
\todo & \textbf{8.4} & \code{pix=1024} & (พิกเซล 0.39 มม.) \\
\end{tabular}
\normalsize

% ===================== SECTION 4: ANALYSIS PROTOCOL =====================
\section{โปรโตคอลการวิเคราะห์ร่วม}
\begin{enumerate}
  \item โหลด \code{seed_N/{control,fracture}/object/SPR.mhd} ทุก seed; คำนวณ
        $\Delta\mathrm{SPR}_N=\mathrm{SPR}^{\text{frac}}_N-\mathrm{SPR}^{\text{ctrl}}_N$
  \item \textbf{a-priori ROI} 16$\times$16px ที่ตำแหน่งร่อง (กลางภาพเมื่อจัด \code{obj_dx/obj_dy}), กำหนดก่อนดูผล;
        $t=\overline{\Delta\mathrm{SPR}}/(s/\sqrt N)$, df$=N-1$
  \item \textbf{background} ROI ขนาดเท่ากัน เลื่อน $\sim18$px --- ต้องได้ null
  \item รายงานเสมอ: $\overline{\Delta\mathrm{SPR}}$, $t$, $p$ (two-sided), $t/p$ ของ background, โฟตอน/พิกเซล
  \item กวาดหลายค่า $\to$ พิจารณา \textbf{multiple-comparison} (เช่น Bonferroni) ก่อนสรุปข้ามเงื่อนไข
  \item เครื่องหมายคาดว่า $<0$ ที่ร่อง --- ค่าบวก ``มีนัยสำคัญ'' = ธงเตือน artifact/misalignment
\end{enumerate}

% ===================== SECTION 5: RUN REFERENCE =====================
\section{อ้างอิงคำสั่งรัน}
\subsection{คู่โมเดล (fracture / matched control)}
\small
\begin{tabular}{c l l}
\rowcolor{softgray} ขนาด & \code{fracture_stl} & \code{control_stl} \\
0.306 & \code{0.306x1000.stl}  & \code{0.306_filled_control.stl} \\
0.43  & \code{0.43x1000.stl}   & \code{0.43x1000_filled_control.stl} \\
0.54  & \code{0.54x1000.stl}   & \code{0.54x1000_filled_control.stl} \\
1.0   & \code{1mmx1000.stl}    & \code{1mmx1000_filled_control.stl} \\
\end{tabular}
\normalsize

\subsection{วิธี A --- job \code{run_seeds} ผ่านแดชบอร์ด (แนะนำ)}
ส่ง job ชนิด \code{run_seeds} ตั้งพารามิเตอร์ตามรายการย่อยของแต่ละ CL (ที่ไม่ระบุ = default):
\begin{run}
type         = run_seeds
control_stl  = 0.54x1000_filled_control.stl
fracture_stl = 0.54x1000.stl
n_seeds      = 10          # 2..30 (สูงสุดต่อ job = 30)
energy_keV   = 80          # CL-3
field_mm     = 40          # CL-4  (0 = เต็มฟิล์ม)
odd          = 400         # CL-5
src_x , src_z              # CL-6 (มุมลำแสง)
pix          = 256         # CL-8
photons      = 2000000     # CL-1
obj_dx = -21.5 , obj_dy = -22.3
mode         = max
\end{run}
ระบบรัน control+fracture ต่อ seed (seed เดียวกัน) แล้วเรียก \code{aggregate-seeds} อัตโนมัติ;
ผลอยู่ที่ \code{runs/remote_jobs/<job_id>/seed_*/}

\subsection{วิธี B --- CLI ท้องถิ่น}
\begin{run}
python -m gate_pelvis run --stl models/0.54x1000.stl --out OUTDIR \
  --energy_keV 80 --field_mm 40 --odd 400 --pix 256 \
  --photons 2000000 --obj_dx -21.5 --obj_dy -22.3 --random_seed 1234567 --mode max
# ทำซ้ำกับ control_stl + เปลี่ยน --random_seed ต่อ seed แล้วรวม:
python -m gate_pelvis aggregate-seeds --root ROOTDIR --n 10 --out ROOTDIR/seed_analysis
\end{run}

\subsection{ช่วงค่าที่ระบบยอมรับ}
\small
\code{energy_keV} 1--1000 (คลินิก 20--150) \quad \code{sod} 50--5000 \quad \code{odd} 50--2000 \quad
\code{pix} 16--2048 \quad \code{field_mm} $\ge0$ \quad \code{src_x/y/z} $\pm5000$ \quad \code{n_seeds} 2--30
\normalsize

% ===================== SECTION 6: PRIORITY =====================
\section{ลำดับความสำคัญที่แนะนำ}\label{sec:priority}
\begin{enumerate}
  \item \textbf{CL-2} (null/false-positive) --- ตรวจความถูกต้อง pipeline ก่อนเชื่อผลอื่น
  \item \textbf{CL-1} (convergence) --- ยืนยัน 0.306mm \& ตรึง seed/photons ที่ใช้ตลอด
  \item \textbf{CL-6} (beam angle) \& \textbf{CL-3} (energy) --- สองการทดสอบกลไก/เชิงสาเหตุที่แข็งแรงสุด
  \item \textbf{CL-7} (ROI/profile) --- อธิบาย plateau (วิเคราะห์ข้อมูลเดิม ต้นทุนต่ำ)
  \item \textbf{CL-4, CL-5, CL-8} --- ปรับให้เหมาะสม (collimation, air-gap, พิกเซล) เพื่อ CNR สูงสุด
\end{enumerate}

% ===================== SECTION 7: LOG TEMPLATE =====================
\section{แบบฟอร์มบันทึกผล}\label{sec:log}
บันทึกทุกการรันหนึ่งแถว:
\renewcommand{\arraystretch}{1.55}
\small
\begin{longtable}{c c l c c c c c}
\rowcolor{theme}\color{white}CL.n & \color{white}job\_id & \color{white}พารามิเตอร์ที่เปลี่ยน &
  \color{white}ขนาด & \color{white}$\Delta$SPR & \color{white}$t$ & \color{white}$p$ & \color{white}bg $p$ \\
\endhead
\ &\ &\ &\ &\ &\ &\ &\ \\ \hline
\rowcolor{rowtint}\ &\ &\ &\ &\ &\ &\ &\ \\ \hline
\ &\ &\ &\ &\ &\ &\ &\ \\ \hline
\rowcolor{rowtint}\ &\ &\ &\ &\ &\ &\ &\ \\ \hline
\ &\ &\ &\ &\ &\ &\ &\ \\ \hline
\rowcolor{rowtint}\ &\ &\ &\ &\ &\ &\ &\ \\ \hline
\ &\ &\ &\ &\ &\ &\ &\ \\ \hline
\rowcolor{rowtint}\ &\ &\ &\ &\ &\ &\ &\ \\ \hline
\ &\ &\ &\ &\ &\ &\ &\ \\ \hline
\bottomrule
\end{longtable}
\normalsize

\vfill
\noindent{\color{theme!45}\rule{\textwidth}{0.4pt}}\\
{\small\color{ink} เอกสารประกอบ repository \code{GATE-Xray-pelvis-SPR}. ทุกการทดลองใช้โมเดลที่มีอยู่
โดยปรับพารามิเตอร์การรันเท่านั้น. ผลฐาน CL-0 = \code{report/report.pdf}.}

\end{document}
